\documentclass{article}
\usepackage[utf8]{inputenc}
\usepackage{amsmath,amssymb}
\usepackage[most]{tcolorbox}
\usepackage{geometry}
\geometry{margin=2.5cm}

\newcommand{\step}[1]{\par\noindent\hangindent=3.5em\hangafter=1%
  \makebox[3.5em][l]{\textbf{#1.}}}
\newcommand{\steptag}[1]{\hfill{\small\textsf{[#1]}}}

\newtcolorbox{proofbox}[1]{
  colback=gray!5, colframe=gray!60,
  fonttitle=\bfseries, title={#1},
  breakable, enhanced,
  left=4pt, right=4pt, top=4pt, bottom=4pt,
  before skip=10pt, after skip=10pt
}

\begin{document}

\begin{proofbox}{Controlled exact-unit C*-rectification: there are univers...}
\step{1} Controlled exact-unit C*-rectification: there are universal C\_rect >= 1 and e\_rect in (0, 1/C\_rect] such that every finite-dimensional epsilon\_X-C*-algebra with 0 <= epsilon\_X <= e\_rect admits, on the same involutive normed space, a bilinear product bold-dot and J = J\textasciicircum{}dagger for which (calX, J, bold-dot, dagger) satisfies EVERY exact-unit epsilon\_r-C*-algebra axiom of def-epsilon-cstar-algebra, including ||J|| = 1, where epsilon\_r = C\_rect*epsilon\_X, and ||J - I\_X|| <= C\_rect*epsilon\_X, ||x bold-dot y - xy|| <= C\_rect*epsilon\_X*||x||*||y||. \steptag{validated} \\[6pt]
\step{1.1} Norming-functional step: for the original self-adjoint approximate unit I\_X with a=||I\_X|| and epsilon\_X<=1/100, there exists a complex-linear functional phi with phi(I\_X)=1, ||phi||<=1/a, and phi(x\textasciicircum{}dagger)=conjugate(phi(x)) for all x. \steptag{validated} \\[6pt]
\step{1.1.1} Finite-dimensional norming extension: for every nonzero u in a finite-dimensional complex normed space there is complex-linear f with ||f||=1 and f(u)=||u||. Indeed on C u put f\_0(lambda u)=lambda||u||; extend G\_0=Re(f\_0) as a real-linear functional dominated by the norm, one real dimension at a time, defining G(m+t v)=G(m)+t c with c chosen between sup\_m(G(m)-||m-v||) and inf\_m(||m+v||-G(m)); the domination of G makes the lower endpoint at most the upper endpoint by the triangle inequality. After finitely many extensions set f(x)=G(x)-iG(i x). Then f is complex-linear and extends f\_0; if theta=arg(f(x)), |f(x)|=Re(e\textasciicircum{}\{-i theta\}f(x))=G(e\textasciicircum{}\{-i theta\}x)<=||x||, while equality holds at u, hence ||f||=1. \steptag{validated} \\[4pt]
\step{1.1.2} Star symmetrization: apply node 1.1.1 to u=I\_X, where a=||I\_X||>=1-epsilon\_X>0, and let f\textasciicircum{}sharp(x)=conjugate(f(x\textasciicircum{}dagger)). The original dagger is a conjugate-linear isometry and I\_X\textasciicircum{}dagger=I\_X by def-epsilon-cstar-algebra, so f\textasciicircum{}sharp is complex-linear, ||f\textasciicircum{}sharp||=1, and f\textasciicircum{}sharp(I\_X)=a. Thus g=(f+f\textasciicircum{}sharp)/2 has ||g||<=1, g(I\_X)=a, and g(x\textasciicircum{}dagger)=conjugate(g(x)); phi=g/a proves node 1.1. \steptag{validated} \\[4pt]
\step{1.2} Explicit rectification step: with phi from node 1.1, define d\_L(y)=y-I\_X y, d\_R(x)=x-x I\_X, h=I\_X I\_X-I\_X, m\_0(x,y)=xy+phi(x)d\_L(y)+phi(y)d\_R(x)+phi(x)phi(y)h, J=I\_X/a, and x bold-dot y=a m\_0(x,y); then bold-dot is bilinear, J=J\textasciicircum{}dagger is its exact two-sided unit with ||J||=1, and the involution is exactly anti-multiplicative for bold-dot. \steptag{validated} \\[6pt]
\step{1.2.1} Unit and normalization algebra: m\_0 is bilinear. Since phi(I\_X)=1, d\_L(I\_X)=d\_R(I\_X)=I\_X-I\_X I\_X=-h, direct substitution gives m\_0(I\_X,y)=y and m\_0(x,I\_X)=x. Hence J=I\_X/a is a two-sided unit for x bold-dot y=a m\_0(x,y); also J\textasciicircum{}dagger=J and ||J||=1 because I\_X\textasciicircum{}dagger=I\_X and a=||I\_X||>0. \steptag{validated} \\[4pt]
\step{1.2.2} Exact star algebra: d\_L(y)\textasciicircum{}dagger=d\_R(y\textasciicircum{}dagger), d\_R(x)\textasciicircum{}dagger=d\_L(x\textasciicircum{}dagger), h\textasciicircum{}dagger=h, and phi(x\textasciicircum{}dagger)=conjugate(phi(x)). Using the original exact identity (xy)\textasciicircum{}dagger=y\textasciicircum{}dagger x\textasciicircum{}dagger from def-epsilon-cstar-algebra term-by-term in the formula for m\_0 yields m\_0(x,y)\textasciicircum{}dagger=m\_0(y\textasciicircum{}dagger,x\textasciicircum{}dagger); multiplication by the real scalar a gives (x bold-dot y)\textasciicircum{}dagger=y\textasciicircum{}dagger bold-dot x\textasciicircum{}dagger. \steptag{validated} \\[6pt]
\step{1.2.2.1} Dependency bridge and explicit adjoint expansion: depend on node 1.1, which supplies the same functional phi used in node 1.2 and proves phi(z\textasciicircum{}dagger)=conjugate(phi(z)) for every z. Since I\_X\textasciicircum{}dagger=I\_X and the original product obeys (uv)\textasciicircum{}dagger=v\textasciicircum{}dagger u\textasciicircum{}dagger, d\_L(y)\textasciicircum{}dagger=(y-I\_X y)\textasciicircum{}dagger=y\textasciicircum{}dagger-y\textasciicircum{}dagger I\_X=d\_R(y\textasciicircum{}dagger), d\_R(x)\textasciicircum{}dagger=(x-x I\_X)\textasciicircum{}dagger=x\textasciicircum{}dagger-I\_X x\textasciicircum{}dagger=d\_L(x\textasciicircum{}dagger), and h\textasciicircum{}dagger=(I\_X I\_X-I\_X)\textasciicircum{}dagger=I\_X I\_X-I\_X=h. Conjugate-linearity of dagger and the scalar identity from node 1.1 therefore give m\_0(x,y)\textasciicircum{}dagger=y\textasciicircum{}dagger x\textasciicircum{}dagger+phi(x\textasciicircum{}dagger)d\_R(y\textasciicircum{}dagger)+phi(y\textasciicircum{}dagger)d\_L(x\textasciicircum{}dagger)+phi(x\textasciicircum{}dagger)phi(y\textasciicircum{}dagger)h=m\_0(y\textasciicircum{}dagger,x\textasciicircum{}dagger), where the last equality only reorders complex scalar factors. Finally a=||I\_X|| is positive real, so (x bold-dot y)\textasciicircum{}dagger=(a m\_0(x,y))\textasciicircum{}dagger=a m\_0(y\textasciicircum{}dagger,x\textasciicircum{}dagger)=y\textasciicircum{}dagger bold-dot x\textasciicircum{}dagger. \steptag{validated} \\[4pt]
\step{1.2.3} Rectification bounds before and after normalization: the approximate-unit axioms give ||d\_L(y)||<=epsilon\_X||y||, ||d\_R(x)||<=epsilon\_X||x||, and ||h||<=epsilon\_X a; with ||phi||<=1/a this implies ||m\_0(x,y)-xy||<=3 epsilon\_X/a ||x||||y||. Since |a-1|<=epsilon\_X and ||xy||<=(1+epsilon\_X)||x||||y||, Delta=a m\_0-m obeys ||Delta(x,y)||<=(3 epsilon\_X+epsilon\_X(1+epsilon\_X))||x||||y||<=5 epsilon\_X||x||||y|| for epsilon\_X<=1/100. This completes node 1.2 including the product-closeness assertion used later. \steptag{validated} \\[4pt]
\step{1.3} Quantitative axiom-transfer step: if Delta(x,y)=x bold-dot y-xy, then ||Delta(x,y)||<=5 epsilon\_X||x||||y||, and the product-norm, associator, C*-lower-bound, conjugate-linear isometric involution, and exact-unit axioms of def-epsilon-cstar-algebra all hold for bold-dot with parameter 100 epsilon\_X. \steptag{validated} \\[6pt]
\step{1.3.1} All nonassociative axioms transfer: by node 1.2.3 put rho=5 epsilon\_X so ||Delta(x,y)||<=rho||x||||y||. Then ||x bold-dot y||<= (1+epsilon\_X+rho)||x||||y||=(1+6 epsilon\_X)||x||||y||<=(1+100 epsilon\_X)||x||||y||. Also ||x\textasciicircum{}dagger bold-dot x||>=||x\textasciicircum{}dagger x||-rho||x|$|\textasciicircum{}2\rangle$=(1-6 epsilon\_X)||x|$|\textasciicircum{}2\rangle$=(1-100 epsilon\_X)||x||\textasciicircum{}2. The unchanged dagger remains a conjugate-linear isometric involution; node 1.2.2 gives its exact product-reversal law, and node 1.2.1 gives the self-adjoint norm-one exact unit. Thus every axiom except the associator estimate holds with epsilon\_r=100 epsilon\_X. \steptag{validated} \\[6pt]
\step{1.3.1.1} Direct exact product-reversal bridge, independent of pending node 1.2.2: use the validated star-symmetrized functional of node 1.1.2, so phi(z\textasciicircum{}dagger)=conjugate(phi(z)). From I\_X\textasciicircum{}dagger=I\_X and the original exact reversal law in def-epsilon-cstar-algebra, d\_L(y)\textasciicircum{}dagger=(y-I\_X y)\textasciicircum{}dagger=y\textasciicircum{}dagger-y\textasciicircum{}dagger I\_X=d\_R(y\textasciicircum{}dagger), d\_R(x)\textasciicircum{}dagger=(x-x I\_X)\textasciicircum{}dagger=x\textasciicircum{}dagger-I\_X x\textasciicircum{}dagger=d\_L(x\textasciicircum{}dagger), and h\textasciicircum{}dagger=(I\_X I\_X-I\_X)\textasciicircum{}dagger=I\_X I\_X-I\_X=h. Therefore, for m\_0(x,y)=xy+phi(x)d\_L(y)+phi(y)d\_R(x)+phi(x)phi(y)h, conjugate-linearity of dagger gives m\_0(x,y)\textasciicircum{}dagger=y\textasciicircum{}dagger x\textasciicircum{}dagger+phi(x\textasciicircum{}dagger)d\_R(y\textasciicircum{}dagger)+phi(y\textasciicircum{}dagger)d\_L(x\textasciicircum{}dagger)+phi(x\textasciicircum{}dagger)phi(y\textasciicircum{}dagger)h=m\_0(y\textasciicircum{}dagger,x\textasciicircum{}dagger), where only commutativity of complex scalars is used in matching the last three terms. Since a=||I\_X|| is a positive real number, (x bold-dot y)\textasciicircum{}dagger=(a m\_0(x,y))\textasciicircum{}dagger=a m\_0(y\textasciicircum{}dagger,x\textasciicircum{}dagger)=y\textasciicircum{}dagger bold-dot x\textasciicircum{}dagger. This proves the missing axiom directly without invoking node 1.2.2. \steptag{validated} \\[4pt]
\step{1.3.1.2} Validated dependency bridge for all remaining nonassociative axioms: import nodes 1.2.1 and 1.2.3. Node 1.2.1 supplies that bold-dot is bilinear and that J is a self-adjoint two-sided exact unit with ||J||=1. Node 1.2.3 supplies ||Delta(x,y)||<=5 epsilon\_X||x||||y||. The original epsilon\_X-C*-axioms give ||xy||<=(1+epsilon\_X)||x||||y|| and ||x\textasciicircum{}dagger x||>=(1-epsilon\_X)||x||\textasciicircum{}2, while dagger is unchanged, conjugate-linear, isometric, and involutive. Hence ||x bold-dot y||<=||xy||+||Delta(x,y)||<=(1+6 epsilon\_X)||x||||y||<=(1+100 epsilon\_X)||x||||y||, and ||x\textasciicircum{}dagger bold-dot x||>=||x\textasciicircum{}dagger x||-||Delta(x\textasciicircum{}dagger,x)||>=(1-6 epsilon\_X)||x|$|\textasciicircum{}2\rangle$=(1-100 epsilon\_X)||x||\textasciicircum{}2, using ||x\textasciicircum{}dagger||=||x||. Together with the exact product-reversal law proved in child 1.3.1.1, these are precisely every def-epsilon-cstar-algebra axiom other than the associator estimate at epsilon\_r=100 epsilon\_X. \steptag{validated} \\[6pt]
\step{1.3.1.2.1} Dependency-complete nonassociative-axiom assembly: depend explicitly on validated nodes 1.2.1, 1.2.3, and 1.3.1.1. Node 1.2.1 gives bilinearity of bold-dot and a self-adjoint two-sided exact unit J with ||J||=1. Node 1.2.3 gives ||Delta(u,v)||<=5 epsilon\_X||u||||v||. Therefore the original product-norm and C*-lower axioms imply ||x bold-dot y||<=||xy||+||Delta(x,y)||<=(1+6 epsilon\_X)||x||||y||<=(1+100 epsilon\_X)||x||||y|| and ||x\textasciicircum{}dagger bold-dot x||>=||x\textasciicircum{}dagger x||-||Delta(x\textasciicircum{}dagger,x)||>=(1-6 epsilon\_X)||x|$|\textasciicircum{}2\rangle$=(1-100 epsilon\_X)||x||\textasciicircum{}2, since ||x\textasciicircum{}dagger||=||x||. The unchanged dagger retains conjugate-linearity, isometry, and involutivity from the original epsilon\_X-C*-algebra, while node 1.3.1.1 supplies the exact missing identity (x bold-dot y)\textasciicircum{}dagger=y\textasciicircum{}dagger bold-dot x\textasciicircum{}dagger. Thus, by the exhaustive axiom list in def-epsilon-cstar-algebra, bilinearity, product norm, conjugate-linear isometric involution, exact product reversal, the C*-lower bound, and the self-adjoint norm-one exact two-sided unit all hold at epsilon\_r=100 epsilon\_X; only the associator estimate remains outside this node. \steptag{validated} \\[4pt]
\step{1.3.2} Associator transfer: write m for the original product, m'=m+Delta for bold-dot, and rho=5 epsilon\_X. Expanding A'=m'(m'(x,y),z)-m'(x,m'(y,z)) gives the original associator plus m(Delta(x,y),z)-m(x,Delta(y,z))+Delta(m'(x,y),z)-Delta(x,m'(y,z)). The original bounds and ||m'(u,v)||<=(1+6 epsilon\_X)||u||||v|| therefore give ||A'||<=[epsilon\_X+2(1+epsilon\_X)rho+2(1+6 epsilon\_X)rho]||x||||y||||z||=[21 epsilon\_X+70 epsilon\_X\textasciicircum{}2]||x||||y||||z||<=100 epsilon\_X||x||||y||||z|| for epsilon\_X<=1/100. Together with node 1.3.1 this proves node 1.3. \steptag{validated} \\[4pt]
\step{1.4} Constant-closing step: choose the universal constants C\_rect=100 and e\_rect=1/100; nodes 1.1-1.3 then give every asserted exact-unit epsilon\_r-C*-algebra axiom with epsilon\_r=C\_rect epsilon\_X and both required closeness bounds, so they establish node 1. \steptag{validated} \\[6pt]
\step{1.4.1} Dependency-gated closure. This step depends on nodes 1.1, 1.2, and 1.3 and may be accepted only after all three are validated. Set C\_rect=100 and e\_rect=1/100. These are universal, C\_rect>=1, and 0<e\_rect=1/C\_rect. For any 0<=epsilon\_X<=e\_rect, node 1.2 supplies on the same involutive normed space the bilinear product bold-dot and the self-adjoint exact unit J=I\_X/a of norm one, together with ||x bold-dot y-xy||<=5 epsilon\_X||x||||y||; node 1.3 supplies every remaining exact-unit epsilon\_r-C*-algebra axiom for epsilon\_r=100 epsilon\_X. Moreover, since a=||I\_X|| and the original approximate-unit axiom gives |a-1|<=epsilon\_X, ||J-I\_X||=||(1/a-1)I\_X||=|1-a|<=epsilon\_X<=100 epsilon\_X. The product bound 5 epsilon\_X||x||||y||<=100 epsilon\_X||x||||y|| gives the other required closeness estimate. Hence all clauses of node 1 hold. \steptag{validated} \\[4pt]
\end{proofbox}

\end{document}
